% Autor szablonu: dr Roman Czapla
% Kompilacja: pdfLaTeX (dwukrotnie, aby zaktualizować spis treści).
% Wymagane logo uczelni: images/logo.png (ścieżka względem pliku głównego).
\documentclass[11pt,a4paper,oneside]{article}
\usepackage[utf8]{inputenc}
\usepackage[T1]{fontenc}
\usepackage[polish,provide=*]{babel}
\usepackage{lmodern,microtype}
\usepackage[margin=25mm,headheight=14pt]{geometry}
\usepackage{xcolor,graphicx,tabularx,booktabs,array,enumitem,fancyhdr,titlesec}
\usepackage[hidelinks]{hyperref}
\definecolor{accent}{HTML}{344B59}
\definecolor{hintcolor}{HTML}{58636A}
\hypersetup{pdfauthor={},pdftitle={Projekt inżynierski - Dokumentacja instalacji i konfiguracji}}
\setlength{\parindent}{0pt}
\setlength{\parskip}{6pt}
\renewcommand{\arraystretch}{1.3}
\setlength{\emergencystretch}{2em}
\setlist[itemize]{leftmargin=*,itemsep=2pt,topsep=3pt}
\titleformat{\section}{\Large\bfseries\color{accent}}{\thesection.}{0.5em}{}
\titleformat{\subsection}{\large\bfseries\color{accent}}{\thesubsection.}{0.5em}{}
\pagestyle{fancy}
\fancyhf{}
\fancyhead[L]{\small Projekt inżynierski}
\fancyhead[R]{\small Instalacja i konfiguracja}
\fancyfoot[C]{\small\thepage}
\renewcommand{\headrulewidth}{0.3pt}
\newcolumntype{Y}{>{\raggedright\arraybackslash}X}

% Podpowiedzi są widoczne podczas wypełniania szablonu.
% Przed oddaniem: uzupełnij pola, zastąp przykłady własną treścią
% i zmień poniższy przełącznik na \podpowiedzifalse.
\newif\ifpodpowiedzi
\podpowiedzitrue
%\podpowiedzifalse
\newcommand{\podpowiedz}[1]{\ifpodpowiedzi\par\begingroup\small\color{hintcolor}\textit{Wskazówka: #1}\par\endgroup\fi}
\newcommand{\pole}[1]{\textnormal{[#1]}}

% Dane dokumentu - uzupełnij w jednym miejscu.
% Ścieżkę logo można zmienić poniżej.
\newcommand{\pliklogo}{images/logo.png}
\newcommand{\logouczelni}{%
  \includegraphics[width=0.8\textwidth,height=40mm,keepaspectratio]{\pliklogo}%
}
\newcommand{\tytulprojektu}{\pole{Tytuł projektu}}
\newcommand{\rokakademicki}{\pole{2026/2027}}
\newcommand{\semestrprojektu}{\pole{zimowy / letni}}
\newcommand{\datadokumentu}{\pole{dd.mm.rrrr}}
\newcommand{\wersjadokumentu}{\pole{1.0}}
% Zespół dwuosobowy tylko za zgodą dyrekcji: zmień na \trzeciaosobafalse.
\newif\iftrzeciaosoba
\trzeciaosobatrue
\newcommand{\studentA}{\pole{Imię i nazwisko osoby 1}}
\newcommand{\studentB}{\pole{Imię i nazwisko osoby 2}}
\newcommand{\studentC}{\pole{Imię i nazwisko osoby 3}}
\newcommand{\albumA}{\pole{numer albumu}}
\newcommand{\albumB}{\pole{numer albumu}}
\newcommand{\albumC}{\pole{numer albumu}}
\newcommand{\prowadzacyA}{\pole{Stopień lub tytuł, imię i nazwisko}}
\newcommand{\prowadzacyB}{\pole{Stopień lub tytuł, imię i nazwisko}}

\begin{document}
\begin{titlepage}
\centering
{\large Uniwersytet Komisji Edukacji Narodowej w Krakowie\par}
{\normalsize Instytut Bezpieczeństwa i Informatyki\par}
\vspace{8mm}
\logouczelni\par
\vspace{10mm}
{\Large\bfseries\color{accent} PROJEKT INŻYNIERSKI\par}
\vspace{3mm}
{\large Dokumentacja instalacji i konfiguracji\par}
\vspace{15mm}
{\LARGE\bfseries\tytulprojektu\par}
\vspace{15mm}
{\small ZESPÓŁ PROJEKTOWY\par}
\vspace{3mm}
\begin{tabular}{ll}
\studentA & Nr albumu: \albumA\\
\studentB & Nr albumu: \albumB
\iftrzeciaosoba\\
\studentC & Nr albumu: \albumC\fi
\end{tabular}
\vspace{12mm}

{\small PROWADZĄCY ZAJĘCIA PROJEKTOWE\par}
\vspace{3mm}
\prowadzacyA\par
\prowadzacyB\par
\vfill
Rok akademicki: \rokakademicki\quad Semestr: \semestrprojektu\par
Data dokumentu: \datadokumentu\quad Wersja: \wersjadokumentu\par
\vspace{6mm}
Kraków\par

\end{titlepage}

\tableofcontents
\podpowiedz{Dokument jest przeznaczony dla osoby przygotowującej środowisko, wdrażającej projekt lub administrującej nim. Jego wymagalność i zakres zależą od charakteru projektu i są uzgadniane z prowadzącymi. Dla animacji lub wizualizacji osobna dokumentacja instalacyjna może nie być potrzebna. Jeśli jest wymagana, opisz rzeczywisty sposób przygotowania środowiska; usuń podsekcje nieadekwatne. Uzupełnij pola, usuń przykłady i wyłącz podpowiedzi przed oddaniem.}

\section{Szczegółowa dokumentacja instalacji i konfiguracji}
\subsection{Zakres i wersja projektu}
\podpowiedz{Podaj wersję projektu, której dotyczy instrukcja, oraz opisany wariant wdrożenia, np. lokalny, serwerowy, kontenerowy lub na urządzeniu. Wskaż, czy instrukcja obejmuje uruchomienie gotowego rezultatu, odtworzenie go z plików źródłowych, czy oba warianty. Opisz wymaganą wiedzę i uprawnienia osoby wykonującej czynności.}
\textbf{Wersja projektu:} \pole{oznaczenie i data.}\par
\textbf{Wariant i zakres:} \pole{opis środowiska i czynności.}\par
\textbf{Wymagane uprawnienia:} \pole{zakres dostępu.}

\newpage
\subsection{Wymagania sprzętowe}
\podpowiedz{Podaj minimalne zasoby potrzebne do działania projektu: procesor, pamięć RAM, miejsce na dysku, GPU, urządzenia i połączenia sieciowe, o ile mają znaczenie. Rozróżnij wymagania do uruchomienia oraz budowania, renderowania lub innych prac, jeśli są różne. Wskazuj wymagania sprawdzone albo wyraźnie oznacz szacunki.}
\begin{tabularx}{\textwidth}{@{}l Y Y@{}}
\toprule
Zasób & Wymaganie & Uwagi\\
\midrule
\pole{zasób} & \pole{minimum} & \pole{zastosowanie lub zalecenie}\\
\pole{zasób} & \pole{minimum} & \pole{zastosowanie lub zalecenie}\\
\bottomrule
\end{tabularx}

\subsection{Środowisko uruchomieniowe i zależności}
\podpowiedz{Wymień system operacyjny, środowiska uruchomieniowe, narzędzia, biblioteki i usługi wraz ze sprawdzonymi wersjami. Uwzględnij bazę danych, kontenery, oprogramowanie specjalistyczne i wtyczki, jeśli występują. Podaj pliki określające zależności oraz sposób ich instalacji. Wskaż wymagany dostęp do sieci lub usług zewnętrznych.}
\begin{tabularx}{\textwidth}{@{}Y l Y@{}}
\toprule
Komponent & Wersja & Rola / sposób uzyskania\\
\midrule
\pole{nazwa} & \pole{wersja} & \pole{opis}\\
\pole{nazwa} & \pole{wersja} & \pole{opis}\\
\bottomrule
\end{tabularx}

\subsection{Przygotowanie plików projektu}
\podpowiedz{Wskaż przekazany pakiet i jego strukturę: kod źródłowy, pliki projektu, zasoby, dane, instalatory lub obrazy kontenerów. Dla projektu programistycznego podaj adres repozytorium oraz identyfikator commitu lub tagu odpowiadającego przekazanemu pakietowi. Dla projektu o innym charakterze wskaż uzgodnione miejsce przechowywania materiałów i oznaczenie wersji plików. Instrukcja powinna umożliwiać pracę z zarchiwizowanym kompletem materiałów; sam odnośnik do repozytorium lub zewnętrznej usługi nie zastępuje przekazania wymaganych plików. Opisz rozpakowanie, rozmieszczenie plików i wymagane przygotowanie danych.}
\pole{Pliki, katalogi i czynności przygotowawcze.}

\subsection{Instalacja krok po kroku}
\podpowiedz{Podaj czynności w kolejności wykonania, z rzeczywistymi poleceniami lub opisem działań w interfejsie. Przy poleceniach wskaż katalog roboczy i wymagane uprawnienia. Opisz budowanie projektu, przygotowanie bazy danych i innych komponentów, jeśli dotyczy. Po istotnych krokach podaj oczekiwany efekt. Poniższe punkty zastąp instrukcją dla własnego projektu.}
\begin{enumerate}
\item \pole{instalacja wymaganych komponentów i sprawdzenie ich wersji.}
\item \pole{przygotowanie projektu lub budowanie z plików źródłowych.}
\item \pole{instalacja rozwiązania i wymaganych danych.}
\end{enumerate}
% Krótkie polecenia: \texttt{polecenie}.
% Wielowierszowe polecenia można umieścić w środowisku verbatim.
% Nie zamieszczaj prawdziwych haseł ani kluczy dostępu.


\subsection{Konfiguracja}
\podpowiedz{Wskaż pliki konfiguracyjne, zmienne środowiskowe i ustawienia usług. Wyjaśnij znaczenie parametrów, wymagane wartości i miejsca ich zmiany. Uwzględnij porty, ścieżki, połączenia, konta i uprawnienia, jeśli występują. Używaj przykładowych wartości; nie wpisuj rzeczywistych haseł, tokenów ani kluczy. Wyjaśnij, jak osoba wdrażająca projekt ma dostarczyć własne dane dostępowe.}
\begin{tabularx}{\textwidth}{@{}Y Y Y@{}}
\toprule
Parametr / lokalizacja & Znaczenie & Wartość / sposób ustawienia\\
\midrule
\pole{parametr} & \pole{opis} & \pole{przykład}\\
\pole{parametr} & \pole{opis} & \pole{przykład}\\
\bottomrule
\end{tabularx}
\podpowiedz{Jeśli AI jest częścią rozwiązania, opisz wymagane modele, usługę lub środowisko lokalne, sposób konfiguracji dostępu, sprawdzoną wersję modelu, o ile można ją ustalić, oraz zależności i limity istotne dla uruchomienia. Samo użycie AI podczas tworzenia projektu nie wymaga takiego opisu tutaj.}

\subsection{Uruchomienie i weryfikacja działania}
\podpowiedz{Opisz pierwsze uruchomienie, zatrzymanie i ponowne uruchomienie rozwiązania. Podaj sposób sprawdzenia gotowości komponentów, lokalizację logów i krótki scenariusz potwierdzający poprawność instalacji. Wskaż dane potrzebne do próby oraz konkretny oczekiwany rezultat. Dla złożonego projektu sprawdź także połączenia między komponentami.}
\textbf{Uruchomienie i zatrzymanie:} \pole{czynności lub polecenia.}\par
\textbf{Próba działania:} \pole{czynności i dane testowe.}\par
\textbf{Oczekiwany rezultat:} \pole{efekt potwierdzający poprawną instalację.}

\subsection{Typowe problemy}
\podpowiedz{Opisz rzeczywiście napotkane lub przewidywalne problemy instalacyjne i konfiguracyjne. Podaj objaw, sposób ustalenia przyczyny oraz działanie naprawcze. Wskaż logi lub polecenia diagnostyczne, jeśli pomagają. Dodaj wiersze według potrzeb.}
\begin{tabularx}{\textwidth}{@{}Y Y Y@{}}
\toprule
Objaw & Przyczyna / diagnostyka & Rozwiązanie\\
\midrule
\pole{objaw} & \pole{opis} & \pole{czynności}\\
\pole{objaw} & \pole{opis} & \pole{czynności}\\
\bottomrule
\end{tabularx}


\subsection{Utrzymanie i usunięcie instalacji}
\podpowiedz{Jeśli dotyczy: opisz kopie zapasowe i odtwarzanie danych, aktualizację projektu, powrót do poprzedniej wersji i usunięcie instalacji. Wskaż, które pliki lub dane należy zachować oraz które czynności mogą spowodować ich utratę. Dostosuj zakres do projektu; nie twórz procedur dla funkcji, których rozwiązanie nie ma.}
\pole{Procedury utrzymania i usunięcia instalacji.}

\subsection{Źródła i materiały pomocnicze}
\podpowiedz{Wskaż wykorzystane instrukcje producentów, dokumentację narzędzi oraz przekazane pliki pomocnicze, np. przykładową konfigurację i skrypty instalacyjne. Podaj nazwy i lokalizacje plików; dla źródeł internetowych - tytuł, adres i datę dostępu. Usuń podsekcję, jeśli nie ma takich materiałów.}
\pole{Wykaz źródeł i materiałów.}
\end{document}
